% texte.tex — f168r, vv. 466–508 (conversion initiale des travaux antérieurs)
\vers[n=466]{ΜΝ.}{Mnésiloque}{τὸ μὲν ὦ γυναῖκες ὀξυθυμεῖσθαι σφόδρα}{Que vous soyez, ô femmes, vivement irritées}
\vers[n=467]{}{}{Εὐριπίδῃ, τοιαῦτ᾽ ἀκουούσας κακά,}{contre Euripide, qui si mal parla de vous,}
\vers[n=468]{}{}{οὐ θαυμάσιόν ἐστ᾽, οὐδ᾽ ἐπιζεῖν τὴν χολήν.}{n'est pas pour étonner, ni que votre bile déborde.}
\vers[n=469]{}{}{καὐτὴ γὰρ ἔγωγ᾽, οὕτως ὀναίμην τῶν τέκνων,}{Car moi-même, — ainsi puissé-je être heureuse dans mes enfants, —}
\vers[n=470]{}{}{μισῶ τὸν ἄνδρ᾽ ἐκεῖνον, εἰ μὴ μαίνομαι.}{je hais cet homme, ou je serais folle.}
\vers[n=471]{}{}{ὅμως δ᾽ ἐν ἀλλήλαισι χρὴ δοῦναι λόγον·}{Néanmoins il convient que nous nous expliquions entre nous,}
\vers[n=472]{}{}{αὐταὶ γάρ ἐσμεν, κοὐδεμί᾽ ἔκφορος λόγου.}{car nous sommes seules, et rien ne sera rapporté de ce qui se dit ici.}
\vers[n=473]{}{}{τί ταῦτ᾽ ἔχουσαι ᾽κεῖνον αἰτιώμεθα}{Pourquoi allons-nous l'accusant}
\vers[n=474]{}{}{βαρέως τε φέρομεν, εἰ δύ᾽ ἡμῶν ἢ τρία}{et sommes-nous indignées de ce que, sachant deux ou trois}
\vers[n=475]{}{}{κακὰ ξυνειδὼς εἶπε δρώσας μυρία;}{de nos méfaits, il les a divulgués, quand nous en commettons d'innombrables ?}
\vers[n=476]{}{}{ἐγὼ γὰρ αὐτὴ πρῶτον, ἵνα μἄλλην λέγω,}{Moi-même toute la première, pour ne rien dire des autres,}
\vers[n=477]{}{}{ξύνοιδ᾽ ἐμαυτῇ πολλά γ᾽· ἀλλ᾽ ἐκεῖνο δ᾽ οὖν}{j'en ai sur la conscience un grand nombre, mais voici}
\vers[n=478]{}{}{δεινότατον, ὅτε νύμφη μὲν ἦ τρεῖς ἡμέρας,}{le plus fort. J'étais mariée de trois jours,}
\vers[n=479]{}{}{ὁ δ᾽ ἀνὴρ παρ᾽ ἐμοὶ καθηῦδεν· ἦν δέ μοι φίλος,}{et mon mari dormait à mes côtés. Or j'avais un amant}
\vers[n=480]{}{}{ὅσπερ με διεκόρησεν οὖσαν ἑπτέτιν.}{qui m'avait déflorée à sept ans.}
\vers[n=481]{}{}{οὖτος πόθῳ μου ᾽κνυεν ἐλθὼν τὴν θύραν·}{Pressé de désirs, il s'en vient gratter à la porte.}
\vers[n=482]{}{}{κᾆτ᾽ εὐθὺς ἔγνων· εἶτα καταβαίνω λάθρᾳ.}{Je comprends sur-le-champ et descends à petit bruit.}
\vers[n=483]{}{}{ὁ δ᾽ ἀνὴρ ἐρωτᾷ `ποῖ σὺ καταβαίνεις;' `ὅποι;}{Mon mari m'interroge : « Où descends-tu? — Où ?}
\vers[n=484]{}{}{στρόφος μ᾽ ἔχει τὴν γαστέρ᾽ ὦνερ κὠδύνη·}{Je suis prise de tranchées, mon homme, et j'ai mal.}
\vers[n=485]{}{}{ἐς τὸν κοπρῶν᾽ οὖν ἔρχομαι.' `βάδιζέ νυν.'}{Aussi vais-je aux latrines. — Va donc ».}
\vers[n=486]{}{}{κᾆθ᾽ ὁ μὲν ἔτριβεν κεδρίδας ἄννηθον σφάκον·}{Et là-dessus il se met à broyer des grains de genévrier, de l'aneth, de la sauge.}
\vers[n=487]{}{}{ἐγὼ δὲ καταχέασα τοῦ στροφέως ὕδωρ}{Pour moi, ayant versé de l'eau sur les gonds,}
\vers[n=488]{}{}{ἐξῆλθον ὡς τὸν μοιχόν· εἶτ᾽ ἠρειδόμην}{j'allai trouver mon galant et subir son assaut}
\vers[n=489]{}{}{παρὰ τὸν Ἀγυιᾶ κύβδ᾽ ἐχομένη τῆς δάφνης.}{à côté de l'Aguieus, courbant l'échine et me tenant au laurier.}
\vers[n=490]{}{}{ταῦτ᾽ οὐδεπώποτ᾽ εἶφ᾽, ὁρᾶτ᾽, Εὐριπίδης·}{Cela au grand jamais, voyez-vous, Euripide n'en a parlé.}
\vers[n=491]{}{}{οὐδ᾽ ὡς ὑπὸ τῶν δούλων τε κὠρεωκόμων}{Non plus qu'il ne dit qu'à défaut d'autres,}
\vers[n=492]{}{}{σποδούμεθ᾽, ἢν μὴ 'χωμεν ἕτερον, οὐ λέγει·}{nous nous faisons hocher par nos esclaves et par des muletiers ;}
\vers[n=493]{}{}{οὐδ᾽ ὡς ὅταν μάλισθ᾽ ὑπό του ληκώμεθα}{ni qu'après avoir paillardé à cœur joie}
\vers[n=494]{}{}{τὴν νύχθ᾽, ἕωθεν σκόροδα διαμασώμεθα,}{la nuit, nous mâchons de l'ail le matin,}
\vers[n=495]{}{}{ἵν᾽ ὀσφρόμενος ἁνὴρ ἀπὸ τείχους εἰσιὼν}{afin que saisi par l'odeur, notre mari, à son retour du rempart,}
\vers[n=496]{}{}{μηδὲν κακὸν δρᾶν ὑποτοπῆται. ταῦθ᾽, ὁρᾷς,}{n'ait aucun soupçon de notre conduite. Cela, voyez-vous,}
\vers[n=497]{}{}{οὐπώποτ᾽ εἶπεν. εἰ δὲ Φαίδραν λοιδορεῖ,}{il n'en a jamais parlé. Que s'il invective Phèdre,}
\vers[n=498]{}{}{ἡμῖν τί τοῦτ᾽ ἔστ᾽; οὐδ᾽ ἐκεῖν᾽ εἴρηκέ πω,}{que nous importe ? Il n'a pas raconté non plus le cas}
\vers[n=499]{}{}{ὡς ἡ γυνὴ δεικνῦσα τἀνδρὶ τοὔγκυκλον}{de cette femme qui, faisant voir son encycle à son mari,}
\vers[n=501]{}{}{τὸν μοιχὸν ἐξέπεμψεν, οὐκ εἴρηκέ πω.}{fit partir son galant ; non, il ne l'a pas raconté.}
\vers[n=502]{}{}{ἑτέραν δ᾽ ἐγᾦδ᾽ ἣ 'φασκεν ὠδίνειν γυνὴ}{Et j'en sais une autre qui prétendit être en travail}
\vers[n=503]{}{}{δέχ᾽ ἡμέρας, ἔως ἐπρίατο παιδίον·}{dix jours durant, jusqu'à ce qu'elle eût acheté un enfant.}
\vers[n=504]{}{}{ὁ δ᾽ ἀνὴρ περιήρχετ᾽ ὠκυτόκι᾽ ὠνούμενος·}{Le mari courait à la ronde faire emplette de drogues pour hâter la délivrance.}
\vers[n=505]{}{}{τὸ δ᾽ εἰσέφερε γραῦς ἐν χύτρᾳ τὸ παιδίον,}{L'enfant fut apporté par une vieille dans une marmite,}
\vers[n=506]{}{}{ἵνα μὴ βοῴη, κηρίῳ βεβυσμένον·}{la bouche bourrée d'un rayon de miel pour qu'il ne criât point.}
\vers[n=507]{}{}{εἶθ᾽ ὡς ἔνευσεν ἡ φέρουσ᾽, εὐθὺς βοᾷ,}{Puis, sur un signe de la porteuse, la femme aussitôt s'écria :}
\vers[n=508]{}{}{`ἄπελθ᾽ ἄπελθ᾽, ἤδη γὰρ ὦνέρ μοι δοκῶ}{« Va-t'en, va-t'en, mon homme, je crois que je vais}
